% !PRINTED 617
\weberpage{617}
%<<CONT
cyklischen Gruppen zu der symmetrischen und der alternirenden Gruppe.

Die aus $s$ hervorgehende Permutation besteht aus einem einzigen Cyklus mit einer ungeraden Gliederzahl $(0, 1, 2 \ldots n-1)$ und gehört daher zu der alternirenden Gruppe. Die Substitution $t$ lässt den Index 0 ungeändert und liefert für die übrigen Ziffern wieder einen einzigen Cyklus. Denn durch $t$ geht 1 in $g$, $g$ in $g^2$, $g^2$ in $g^3$ etc. über und da die Potenzen $1, g, g^2 \ldots g^{n-2}$, von der Reihenfolge abgesehen, mit den Ziffern $1, 2 \ldots n-1$ übereinstimmen, so entspricht $t$ dem Cyklus $(1, g, g^2 \ldots g^{n-2})$, der aus einer geraden Gliederzahl besteht und folglich zu den Permutationen zweiter Art gehört, also nicht in der alternirenden Gruppe enthalten ist. Dagegen ist $t^2$ wieder darin enthalten.

Daraus ergiebt sich das Resultat:

\gesperrt{IV. Die volle metacyklische Gruppe ist kein Theiler der alternirenden Gruppe. Der grösste gemeinschaftliche Theiler beider Gruppen ist die halbmetacyklische Gruppe.}

Jeder transitive Theiler der metacyklischen Gruppe ist selbst metacyklisch.

Man kann der Bedingung für die metacyklischen Gleichungen verschiedene Formen geben, die sich aus dem Bisherigen ableiten lassen.

Die volle lineare Gruppe ist als Theiler vom Index $\nu = 1 . 2 . 3 \ldots (n-2)$ in der symmetrischen Gruppe enthalten. Eine zu der vollen linearen Gruppe gehörige Function $y$ der $n$ Variablen $x_0, x_1 \ldots x_{n-1}$, die wir eine \gesperrt{metacyklische Function} nennen können, genügt daher einer Resolvente $F(y)=0$ vom Grade $\nu$, deren Coëfficienten symmetrische Functionen der $x$ sind.

Substituirt man nun für $x$ die Wurzeln einer metacyklischen Gleichung $f(x)=0$, so wird $y$ rational.

Wenn umgekehrt die Function $y$ durch die Substitution der Wurzeln einer irreduciblen Gleichung $f(x)=0$ für die $x$ rational wird, während $F'(y)$ von Null verschieden bleibt, so ist $f(x)$ metacyklisch; denn dann ist die Gruppe von $f(x)=0$ gewiss ein Theiler der vollen metacyklischen Gruppe, und also, da $f(x)$ irreducibel ist, selbst metacyklisch. Es genügt aber auch für die algebraische Auflösbarkeit von $f(x)=0$, wenn die Resolvente $F(y)=0$ nur überhaupt eine rationale Wurzel hat, die
%CONT>>
